\subsection{哪些结论成立，哪些预期没有泛化}

\paragraph{总体均值有改善，不能据此宣称处处更优。}
冻结的距离域交叉验证方法在360例主测试上的平均误差为4.189米，相比WLS的6.532米降低35.87\%。它比固定恒速稳健基线的均值低10.38\%，但CEP90为9.966米，略高于后者的9.853米。相对旧BIC门控，均值只减少0.158米，配对95\%区间为$[-0.091,0.425]$米，包含零；本轮不能确认其稳定优于BIC。一个总体均值无法同时概括静态、动态与尾部风险。

\paragraph{静止非劣条件没有从验证集推广到测试集。}
验证集静态正常误差为0.902米，低于同批WLS的0.923米；但独立测试中变为1.419米，WLS为1.031米，退化37.63\%。因此，即使把开发判据放宽为10\%容忍度，测试结果仍未满足。更不能把总体35.87\%的均值改善写成完成原题的“标准场景不劣化”或“移动检测精度提高20\%”。原题的后一指标没有得到定义，也没有真实工业对照数据。

\paragraph{门控节省了静态过拟合的代价，也丢失了一部分运动收益。}
固定恒速消融的主测试均值为4.674米，高于所选方案，但其压力测试均值为8.814米，低于所选方案的10.176米。所选方案在稀疏场景的平均误差为11.121米，固定恒速基线为5.976米。这里有可核查的实现机制：三折划分下，8份观测产生不足6条的训练折，运动模型因预设可用性门槛被排除。它不是从数据中“发现目标静止”；这种保守回退需要作为系统失效方式报告。

\paragraph{计算成本提高，稳健性仍有明确边界。}
主测试上所选方法平均消耗71.3个工作单位，旧BIC为47.8、恒速稳健为27.2；本机记录的平均用时约为14.9、7.6和4.9毫秒。工作单位在新旧求解器之间的计数口径有差异，壁钟时间也依赖实现和设备，所以这些数字只说明本次运行的精度与成本取舍。最终配置没有使用截尾或子集初始化，不能把本轮效果归因于这些虽然已实现、但没有进入冻结方案的组件。

\paragraph{不确定性层必须报告实测欠覆盖。}
名义90\%区域在独立同分布混合样本上仅覆盖274/320例，即85.625\%；名义95\%区域覆盖303/320例，即94.6875\%。移动异常子群的相应覆盖为59.76\%和87.80\%，静态两个子群则均为100\%。这展示了总体校准与子群均衡之间的差别。压力分布上的覆盖进一步降至74.0\%和85.5\%，尤其稀疏场景分别只有37.5\%和60.0\%。

共形定理对随机校准样本与新案例联合取概率，不承诺任一已实现校准集下的有限测试比例必定超过名义水平。本次90\%经验覆盖没有达标，这一事实保留；既不能将它掩盖为“90\%覆盖已获验证”，也不能仅凭一个有限样本比例否定带有条件的边际覆盖定理。半径10米的90\%区域拒绝策略在同分布测试中接受71.56\%案例，接受后的覆盖只有83.84\%，再次说明按半径拒绝不会自动保留条件覆盖。

上述结果支持一个有限结论：本研究得到一个对常规混合场景有平均收益的可审计方案，并定位了静止泛化、稀疏门控、尾部风险与子群校准方面的失败。开放问题仍然存在，失败结果本身为下一轮建模提供了更明确的约束。
